\section{Заключение}
\label{sec:Chapter5} \index{Chapter5}

%Здесь надо перечислить все результаты, полученные в ходе работы. Из текста
%должно быть понятно, в какой мере решена поставленная задача.

В ходе выполнения дипломной работы были получены следующие результаты:
\begin{enumerate}
    \item Изучены существующие подходы к снижению влияния NUMA-эффекта и выявлены их ограничения.

    Проведённый обзор показал, что существующие методы не всегда позволяют эффективно устранять последствия удалённых обращений к памяти, особенно для больших cross-NUMA нагрузок, что обосновывает целесообразность разработки механизма репликации таблиц трансляций и данных пользовательских приложений.

    \item Спроектирован механизм репликации таблиц трансляций и данных пользовательских приложений для NUMA-систем.

    Представлена возможность независимого управления репликацией таблиц трансляций и репликацией пользовательских данных, а также реализована поддержка нескольких политик репликации, ориентированных на различные сценарии эксплуатации.

    При проектировании механизма обеспечена его прозрачность: его использование не требует модификации программ, их бинарных файлов или процесса сборки, а также применения специализированных средств управления памятью. Управление механизмом осуществляется через \textit{control groups}\cite{cgroup} путём настройки политик репликации для соответствующей группы процессов. При этом функционирование приложений сохраняется в штатном режиме.

    Кроме того, при проектировании были рассмотрены особенности интеграции механизма репликации с функциональностями \textit{Transparent Huge Pages}\cite{THP} и \textit{HugeTLB}\cite{hugetlb}, что создаёт основу для расширения механизма поддержкой \textit{больших страниц} в дальнейшем.

    \item Спроектированный механизм реализован в ядре дистрибутива openEuler 25.03, основанного на Linux 6.6, для архитектуры arm64. Поддержаны 4K и 64K размеры страниц.

    \item Выполнена верификация разработанного механизма, нарушений корректности функционирования не выявлено.

    Все использованные в работе тестовые приложения успешно отработали при различных комбинациях политик репликации без возникновения сбоев ядра и нарушений корректности выполнения программ. Дополнительно модифицированное ядро было протестировано с использованием набора тестов \textit{Linux Test Project}\cite{ltp}. По результатам тестирования не было выявлено ни отказов ядра, ни деградации результатов тестов по сравнению с исходной версией ядра.

    \item Проведена экспериментальная оценка производительности механизма на наборе синтетических и прикладных тестов.

    Полученные результаты показали, что в ряде сценариев репликация таблиц трансляций и данных позволяет снизить влияние NUMA-эффекта и повысить производительность приложений.

    \item Рассмотрены вопросы автоматизации управления механизмом репликации.

    В текущей реализации в качестве политики по умолчанию для всех процессов используется комбинация (\texttt{Minimal}, \texttt{Lazy}), представляющая собой компромисс между полным отсутствием репликации и наиболее агрессивными режимами её применения.

    Построение более точной модели автоматического выбора политик репликации требует отдельного исследования. Такая модель должна учитывать характеристики выполняемых приложений, выявлять признаки, определяющие эффективность репликации для различных типов нагрузок. Для её практического применения также требуется разработка механизмов сбора и анализ соответствующих метрик во время выполнения программ. Решение данной задачи представляет собой перспективное направление дальнейших исследований.
\end{enumerate}
